%=============================================================================
% Comparativa Técnica de Sistemas de Imagen Médica
% Curso: Imagenología Médica — B.Sc. Ingeniería Biomédica, UNMSM
% Autora: Joselyn Romero Avila — 2025
%=============================================================================

\documentclass[11pt, a4paper]{article}

% ── Paquetes ──────────────────────────────────────────────────────────────────
\usepackage[T1]{fontenc}
\usepackage[utf8]{inputenc}
\usepackage{babel}
\usepackage[margin=2.2cm, top=2.5cm, bottom=2.5cm]{geometry}


\usepackage{xcolor}
\usepackage{booktabs}
\usepackage{tabularx}
\usepackage{longtable}
\usepackage{array}
\usepackage{multirow}
\usepackage{colortbl}
\usepackage{enumitem}
\usepackage{hyperref}
\usepackage{fancyhdr}
\usepackage{titlesec}
\usepackage{parskip}
\usepackage{caption}
\usepackage{float}
\usepackage{amsmath}
\usepackage{amssymb}
\usepackage{graphicx}
\usepackage{tcolorbox}
\tcbuselibrary{skins, breakable}

% ── Colores ───────────────────────────────────────────────────────────────────
\definecolor{navy}{HTML}{0A2540}
\definecolor{teal}{HTML}{007B8A}
\definecolor{teallt}{HTML}{E6F4F6}
\definecolor{accent}{HTML}{1A5276}
\definecolor{grayd}{HTML}{2C3E50}
\definecolor{graym}{HTML}{7F8C8D}
\definecolor{grayl}{HTML}{F4F6F8}
\definecolor{graybrd}{HTML}{D5DBDB}
\definecolor{bestgreen}{HTML}{D5F5E3}
\definecolor{rowalt}{HTML}{EBF5FB}

% ── Hipervínculos ─────────────────────────────────────────────────────────────
\hypersetup{
  colorlinks=true,
  linkcolor=accent,
  urlcolor=teal,
  hidelinks
}

% ── Tipografía de secciones ───────────────────────────────────────────────────
\titleformat{\section}
  {\color{navy}\Large\bfseries}
  {\thesection.}{0.6em}{}
  [\vspace{1pt}{\color{teal}\titlerule[1.5pt]}\vspace{4pt}]

\titleformat{\subsection}
  {\color{accent}\normalsize\bfseries}
  {}{0em}{}

\titlespacing{\section}{0pt}{16pt}{6pt}
\titlespacing{\subsection}{0pt}{10pt}{4pt}

% ── Header / Footer ───────────────────────────────────────────────────────────
\pagestyle{fancy}
\fancyhf{}
\renewcommand{\headrulewidth}{0pt}
\fancyhead[L]{\small\color{graym}\textit{Comparativa de Sistemas de Imagen Médica}}
\fancyhead[R]{\small\color{graym}\textit{Imagenología Médica · UNMSM · 2025}}
\fancyfoot[C]{\small\color{graym}\thepage}
\fancyfoot[L]{\small\color{graym}\textit{Joselyn Romero Avila · Ing. Biomédica}}
\setlength{\headheight}{14pt}

% ── Tabla de parámetros: columnas ─────────────────────────────────────────────
\newcolumntype{P}[1]{>{\raggedright\arraybackslash}p{#1}}
\newcolumntype{C}[1]{>{\centering\arraybackslash}p{#1}}

% ── Comando: fila de categoría ────────────────────────────────────────────────
\newcommand{\catrow}[1]{%
  \rowcolor{teallt}
  \multicolumn{1}{l}{\textbf{\small\color{accent}#1}} &
  \multicolumn{1}{c}{} &
  \multicolumn{1}{c}{} &
  \multicolumn{1}{c}{} &
  \multicolumn{1}{c}{} &
  \multicolumn{1}{c}{} \\[-2pt]
}

% ── Comando: caja de principio físico ────────────────────────────────────────
\newtcolorbox{principiobox}{
  colback=teallt,
  colframe=teal,
  boxrule=0.8pt,
  arc=4pt,
  left=8pt, right=8pt, top=6pt, bottom=6pt,
  fonttitle=\bfseries\small\color{navy},
  title=Principio físico,
}

% ── Comando: caja de aplicaciones clínicas ───────────────────────────────────
\newtcolorbox{clinicbox}{
  colback=grayl,
  colframe=graybrd,
  boxrule=0.5pt,
  arc=4pt,
  left=8pt, right=8pt, top=6pt, bottom=6pt,
  fonttitle=\bfseries\small\color{navy},
  title=Aplicaciones clínicas principales,
  breakable,
}

%=============================================================================
\begin{document}
\pagenumbering{arabic}

%-----------------------------------------------------------------------------
% PORTADA
%-----------------------------------------------------------------------------
\thispagestyle{empty}

\vspace*{1.5cm}

{\noindent\color{teal}\rule{\linewidth}{2pt}}
\vspace{0.3cm}

{\noindent\fontsize{9}{11}\selectfont\color{graym}\textsc{Reporte Técnico — Proyecto Final de Curso}}

\vspace{0.4cm}

{\noindent\fontsize{24}{30}\selectfont\bfseries\color{navy}
Comparativa Técnica de Sistemas\\[4pt]
de Imagen Médica}

\vspace{0.5cm}

{\noindent\fontsize{12}{16}\selectfont\color{grayd}
Análisis de parámetros técnicos y aplicaciones clínicas:\\
MRI · CT · Ultrasonido · Radiografía · PET/CT · Fluoroscopía}

\vspace{0.8cm}
{\noindent\color{graybrd}\rule{\linewidth}{0.5pt}}
\vspace{0.5cm}

\begin{tabular}{@{}ll}
  \textbf{\color{accent}Asignatura:}  & Imagenología Médica \\[3pt]
  \textbf{\color{accent}Programa:}    & B.Sc. Ingeniería Biomédica — UNMSM \\[3pt]
  \textbf{\color{accent}Autora:}      & Joselyn Romero Avila \\[3pt]
  \textbf{\color{accent}Año:}         & 2025 \\[3pt]
  \textbf{\color{accent}Fuentes:}     & Datasheets técnicos públicos de fabricantes (2024--2025) \\
\end{tabular}

\vspace{0.8cm}
{\noindent\color{graybrd}\rule{\linewidth}{0.5pt}}
\vspace{0.5cm}

\noindent\textbf{\color{navy}Resumen}

\vspace{0.3cm}
\noindent Este reporte presenta una comparativa técnica de los principales sistemas de imagen médica
utilizados en entornos clínicos y hospitalarios. Para cada modalidad se describen los principios
físicos de operación, los parámetros técnicos clave y las aplicaciones clínicas prioritarias.
El análisis comparativo por tipo de equipo evalúa parámetros de rendimiento, seguridad,
infraestructura requerida y rango de precio referencial, con el objetivo de orientar decisiones
de adquisición en el sistema de salud.

\vspace{0.8cm}

\noindent\textbf{\color{navy}Índice de contenidos}

\vspace{0.3cm}
\begin{enumerate}[leftmargin=*, label=\textbf{\arabic*.}, itemsep=4pt]
  \item Resonancia Magnética (MRI) \dotfill \textit{1.5T y 3T}
  \item Tomografía Computada (CT) \dotfill \textit{64, 128 y 256 cortes}
  \item Ultrasonido (US) \dotfill \textit{Portátil, mid-range, premium, POCUS}
  \item Radiografía (RX) \dotfill \textit{CR, DR directo, móvil, mamógrafo}
  \item PET/CT \dotfill \textit{Estándar, TOF, digital, PET/MR}
  \item Fluoroscopía e Intervencionismo \dotfill \textit{General, arco en C, hemodinámica, sala híbrida}
  \item Notas sobre especificaciones \dotfill
\end{enumerate}

\vspace{0.5cm}
{\noindent\color{teal}\rule{\linewidth}{2pt}}

\newpage

%=============================================================================
\section{Resonancia Magnética (MRI --- Magnetic Resonance Imaging)}
%=============================================================================

\begin{principiobox}
Excitación de núcleos de hidrógeno ($^1$H) por campo magnético estático ($B_0$) y pulsos de
radiofrecuencia (RF). La adquisición se realiza en el espacio-$k$ y la imagen se obtiene
mediante la transformada de Fourier inversa. \textbf{Sin radiación ionizante.}
\end{principiobox}

\vspace{6pt}
La resonancia magnética utiliza un campo magnético estático ($B_0$), gradientes espaciales y
pulsos de radiofrecuencia para generar imágenes de alta resolución en tejidos blandos.
Los parámetros clave para evaluación técnica incluyen la intensidad del campo (T),
la amplitud de los gradientes (mT/m), el SAR (\textit{Specific Absorption Rate}) y la
cobertura de canales RF.

\subsection{Comparativa de parámetros técnicos}

{\footnotesize $\bigstar$ = valor más favorable en esa categoría de campo magnético}

\vspace{4pt}
\begin{center}
\renewcommand{\arraystretch}{1.25}
\begin{tabular}{P{4.2cm} C{2cm} C{2cm} C{2cm} C{2cm} C{2cm}}
\toprule
\rowcolor{navy}
\textbf{\color{white}Parámetro} &
\textbf{\color{white}\shortstack{Marca A\\\\(1.5T)}} &
\textbf{\color{white}\shortstack{Marca B\\\\(1.5T)}} &
\textbf{\color{white}\shortstack{Marca C\\\\(1.5T)}} &
\textbf{\color{white}\shortstack{Marca A\\\\(3T)}} &
\textbf{\color{white}\shortstack{Marca B\\\\(3T)}} \\
\midrule
\rowcolor{teallt}\small\textbf{Core} & & & & & \\
Campo magnético (T)        & 1.5 & 1.5 & 1.5 & 3.0 & 3.0 \\
\rowcolor{rowalt}
Gradiente máx. (mT/m)      & \cellcolor{bestgreen}$\bigstar$ 45 & 44 & 45 & 60 & \cellcolor{bestgreen}$\bigstar$ 80 \\
Slew rate (T/m/s)          & 200 & 200 & 200 & 200 & \cellcolor{bestgreen}$\bigstar$ 333 \\
\rowcolor{teallt}\small\textbf{Paciente} & & & & & \\
Bore (cm)                  & 70 & 70 & \cellcolor{bestgreen}$\bigstar$ 71 & 70 & 70 \\
\rowcolor{rowalt}
Peso máx. paciente (kg)    & \cellcolor{bestgreen}$\bigstar$ 250 & 227 & \cellcolor{bestgreen}$\bigstar$ 250 & \cellcolor{bestgreen}$\bigstar$ 250 & 227 \\
\rowcolor{teallt}\small\textbf{Imagen} & & & & & \\
FOV máx. (cm)              & 50 & 50 & \cellcolor{bestgreen}$\bigstar$ 53 & 50 & 50 \\
\rowcolor{rowalt}
Canales RF máx.            & \cellcolor{bestgreen}$\bigstar$ 204 & 128 & 128 & \cellcolor{bestgreen}$\bigstar$ 204 & 128 \\
\rowcolor{teallt}\small\textbf{Seguridad} & & & & & \\
SAR límite (W/kg)          & 4.0 & 4.0 & 4.0 & 4.0 & 4.0 \\
\rowcolor{rowalt}
Nivel de ruido (dB)        & \cellcolor{bestgreen}$\bigstar$ 72 & 74 & 73 & \cellcolor{bestgreen}$\bigstar$ 78 & 79 \\
\rowcolor{teallt}\small\textbf{Operación} & & & & & \\
Helio boil-off (L/año)     & \cellcolor{bestgreen}$\bigstar$ 0 & \cellcolor{bestgreen}$\bigstar$ 0 & \cellcolor{bestgreen}$\bigstar$ 0 & \cellcolor{bestgreen}$\bigstar$ 0 & \cellcolor{bestgreen}$\bigstar$ 0 \\
\rowcolor{rowalt}
Consumo eléctrico (kVA)    & \cellcolor{bestgreen}$\bigstar$ 35 & 40 & 38 & \cellcolor{bestgreen}$\bigstar$ 55 & 60 \\
\rowcolor{teallt}\small\textbf{IT / Conectividad} & & & & & \\
DICOM / HL7                & Sí & Sí & Sí & Sí & Sí \\
\rowcolor{rowalt}
IA integrada               & Sí & Sí & Sí & Sí & Sí \\
\rowcolor{teallt}\small\textbf{Comercial} & & & & & \\
Precio ref. (USD)          & 850K--1.2M & 800K--1.1M & 820K--1.15M & 1.4M--2.2M & 1.35M--2.1M \\
\bottomrule
\end{tabular}
\end{center}

\vspace{4pt}
\begin{itemize}[leftmargin=*, itemsep=2pt, topsep=2pt]
  \footnotesize
  \item \textbf{SAR} (\textit{Specific Absorption Rate}): límite regulatorio IEC 60601-2-33, máx. 4~W/kg cuerpo entero
  \item \textbf{Helio boil-off = 0} en equipos modernos gracias a tecnología de imán sellado sin ventilación activa
  \item \textbf{Precio referencial FOB} — varía por configuración, accesorios e infraestructura de sala
\end{itemize}

\begin{clinicbox}
\begin{tabular}{@{}P{3cm} P{10.5cm}@{}}
\textbf{Neurología}     & Secuencias FLAIR, DWI, SWI --- alta sensibilidad en tejido cerebral \\[2pt]
\textbf{Oncología}      & Perfusión, espectroscopía, DWI --- estadificación y seguimiento tumoral \\[2pt]
\textbf{Cardiología}    & Cine MRI, mapeo T1/T2 --- función ventricular y viabilidad miocárdica \\[2pt]
\textbf{MSK}            & Cartílago, ligamentos, tendones --- sin radiación ionizante \\[2pt]
\textbf{Abdomen/Pelvis} & MRCP, angio-MR, secuencias dinámicas con contraste \\
\end{tabular}
\end{clinicbox}

\newpage

%=============================================================================
\section{Tomografía Computada (CT --- Computed Tomography)}
%=============================================================================

\begin{principiobox}
Atenuación diferencial de rayos X según densidad tisular. Reconstrucción volumétrica
mediante retroproyección filtrada o algoritmos iterativos/IA.
Los detectores modernos permiten adquisición helicoidal con tiempos de rotación $< 0.3$~s.
\end{principiobox}

\vspace{6pt}
La reconstrucción iterativa (ASIR, ADMIRE, iDose) y los algoritmos de \textit{deep learning}
(AiCE, TrueFidelity) permiten reducir la dosis de radiación hasta un 80\% respecto a la TC
convencional. La CT espectral (\textit{dual energy}) permite caracterización tisular adicional
mediante análisis por coeficiente de atenuación.

\subsection{Comparativa de parámetros técnicos}

\begin{center}
\renewcommand{\arraystretch}{1.25}
\begin{tabular}{P{4.2cm} C{2cm} C{2cm} C{2cm} C{2cm} C{2cm}}
\toprule
\rowcolor{navy}
\textbf{\color{white}Parámetro} &
\textbf{\color{white}CT 64c (A)} &
\textbf{\color{white}CT 64c (B)} &
\textbf{\color{white}CT 64c (C)} &
\textbf{\color{white}CT 128c (A)} &
\textbf{\color{white}CT 256c (A)} \\
\midrule
\rowcolor{teallt}\small\textbf{Core} & & & & & \\
Cortes por rotación        & 64 & 64 & 64 & 128 & \cellcolor{bestgreen}$\bigstar$ 256 \\
\rowcolor{rowalt}
Velocidad rotación (s)     & 0.50 & 0.50 & 0.50 & 0.33 & \cellcolor{bestgreen}$\bigstar$ 0.28 \\
Cobertura detector (mm)    & 38.4 & 40 & 40 & 76.8 & \cellcolor{bestgreen}$\bigstar$ 141 \\
\rowcolor{rowalt}
Potencia tubo (kW)         & 40 & \cellcolor{bestgreen}$\bigstar$ 50 & 40 & 60 & \cellcolor{bestgreen}$\bigstar$ 100 \\
Rango kVp                  & 80--130 & 80--140 & 80--140 & 70--150 & 70--150 \\
\rowcolor{teallt}\small\textbf{Imagen} & & & & & \\
\rowcolor{rowalt}
Res. espacial (lp/cm)      & 14 & 14 & 14 & 15 & \cellcolor{bestgreen}$\bigstar$ 16 \\
Res. temporal cardíaca (ms)& 165 & 175 & 175 & 83 & \cellcolor{bestgreen}$\bigstar$ 75 \\
\rowcolor{teallt}\small\textbf{Paciente} & & & & & \\
\rowcolor{rowalt}
Bore (cm)                  & \cellcolor{bestgreen}$\bigstar$ 78 & 70 & \cellcolor{bestgreen}$\bigstar$ 78 & \cellcolor{bestgreen}$\bigstar$ 78 & \cellcolor{bestgreen}$\bigstar$ 78 \\
Peso máx. paciente (kg)    & \cellcolor{bestgreen}$\bigstar$ 250 & 227 & \cellcolor{bestgreen}$\bigstar$ 250 & \cellcolor{bestgreen}$\bigstar$ 300 & \cellcolor{bestgreen}$\bigstar$ 300 \\
\rowcolor{teallt}\small\textbf{Seguridad} & & & & & \\
\rowcolor{rowalt}
Reducción de dosis (\%)    & 60 & 55 & 58 & 70 & \cellcolor{bestgreen}$\bigstar$ 80 \\
\rowcolor{teallt}\small\textbf{Operación} & & & & & \\
Consumo eléctrico (kVA)    & \cellcolor{bestgreen}$\bigstar$ 32 & 35 & \cellcolor{bestgreen}$\bigstar$ 32 & 50 & 80 \\
\rowcolor{teallt}\small\textbf{Avanzado} & & & & & \\
\rowcolor{rowalt}
CT Espectral / Dual Energy & No & No & No & Opcional & \cellcolor{bestgreen}$\bigstar$ Sí \\
IA en reconstrucción       & Sí & Sí & Sí & Sí & Sí \\
\rowcolor{teallt}\small\textbf{Comercial} & & & & & \\
\rowcolor{rowalt}
Precio ref. (USD)          & 280K--420K & 270K--410K & 275K--415K & 450K--680K & 700K--950K \\
\bottomrule
\end{tabular}
\end{center}

\begin{itemize}[leftmargin=*, itemsep=2pt, topsep=4pt]
  \footnotesize
  \item Reducción de dosis referida a técnica convencional sin algoritmo iterativo
  \item \textbf{CT Espectral}: permite caracterización tisular por coeficiente de atenuación en dos energías
  \item Precio referencial FOB — varía por configuración y mercado de destino
\end{itemize}

\begin{clinicbox}
\begin{tabular}{@{}P{3.5cm} P{10cm}@{}}
\textbf{Urgencias y Trauma}  & Protocolo whole-body, alta velocidad --- $<$10~s para tórax-abdomen-pelvis \\[2pt]
\textbf{Cardiología}         & Angio coronaria, score de calcio --- requiere CT $\geq$ 64c, res. temporal $<$ 100~ms \\[2pt]
\textbf{Oncología}           & Estadificación, seguimiento RECIST, guía de biopsias \\[2pt]
\textbf{Neurología}          & ACV isquémico/hemorrágico, perfusión cerebral, angio-TC \\[2pt]
\textbf{Tórax / Pulmón}      & Nódulos pulmonares, HRCT para intersticio, embolismo pulmonar \\
\end{tabular}
\end{clinicbox}

\newpage

%=============================================================================
\section{Ultrasonido / Ecografía (US --- Ultrasonography)}
%=============================================================================

\begin{principiobox}
Emisión y recepción de ondas acústicas de alta frecuencia (2--18~MHz). La imagen se forma
por reflexión en interfaces tisulares con diferente impedancia acústica ($Z = \rho \cdot c$).
\textbf{Sin radiación ionizante. Modalidad en tiempo real.}
\end{principiobox}

\vspace{6pt}
Los transductores modernos emplean matrices de cristales piezoeléctricos con formación
digital de haz (\textit{beamforming}). Frecuencias altas (10--18~MHz) ofrecen mayor
resolución pero menor penetración; frecuencias bajas (2--5~MHz) permiten mayor profundidad.
El Doppler color y pulsado evalúan flujo vascular en tiempo real.

\subsection{Comparativa de parámetros técnicos}

\begin{center}
\renewcommand{\arraystretch}{1.25}
\begin{tabular}{P{4.2cm} C{2.5cm} C{2.5cm} C{2.5cm} C{2.5cm}}
\toprule
\rowcolor{navy}
\textbf{\color{white}Parámetro} &
\textbf{\color{white}Portátil} &
\textbf{\color{white}Mid-Range} &
\textbf{\color{white}Premium} &
\textbf{\color{white}POCUS} \\
\midrule
Rango frecuencia (MHz)     & 2--10 & 1--18 & \cellcolor{bestgreen}$\bigstar$ 1--22 & 3--12 \\
\rowcolor{rowalt}
Transductores disp.        & 3--5 & 8--12 & \cellcolor{bestgreen}$\bigstar$ 15+ & 2--3 \\
Doppler color              & Sí & Sí & Sí & Sí \\
\rowcolor{rowalt}
Elastografía               & No & Básica & \cellcolor{bestgreen}$\bigstar$ Avanzada & No \\
3D/4D tiempo real          & No & Sí & \cellcolor{bestgreen}$\bigstar$ Sí & No \\
\rowcolor{rowalt}
Contraste (CEUS)           & No & Opcional & \cellcolor{bestgreen}$\bigstar$ Sí & No \\
Profundidad máx. (cm)      & 25 & \cellcolor{bestgreen}$\bigstar$ 35 & \cellcolor{bestgreen}$\bigstar$ 35 & 20 \\
\rowcolor{rowalt}
Pantalla                   & 15 pulg. & 19 pulg. & \cellcolor{bestgreen}$\bigstar$ 23 pulg. & \cellcolor{bestgreen}$\bigstar$ Tablet \\
Peso equipo (kg)           & \cellcolor{bestgreen}$\bigstar$ 8 & 15 & 25 & \cellcolor{bestgreen}$\bigstar$ 0.3 \\
\rowcolor{rowalt}
Batería autónoma           & Sí & No & No & \cellcolor{bestgreen}$\bigstar$ Sí \\
IA integrada               & Básica & Sí & \cellcolor{bestgreen}$\bigstar$ Avanzada & Básica \\
\rowcolor{rowalt}
Precio ref. (USD)          & 30K--60K & 60K--120K & 120K--250K & 2K--8K \\
\bottomrule
\end{tabular}
\end{center}

\begin{itemize}[leftmargin=*, itemsep=2pt, topsep=4pt]
  \footnotesize
  \item \textbf{POCUS} (\textit{Point-of-Care Ultrasound}): dispositivos ultraminiaturizados para urgencias y UCI
  \item \textbf{CEUS} (\textit{Contrast-Enhanced Ultrasound}): microburbujas para caracterización hepática y vascular
  \item \textbf{Elastografía}: cuantificación de rigidez tisular --- hígado (fibrosis), tiroides, mama
\end{itemize}

\begin{clinicbox}
\begin{tabular}{@{}P{3cm} P{10.5cm}@{}}
\textbf{Abdomen}       & Hígado, vesícula, páncreas, riñones --- primera línea de imagen abdominal \\[2pt]
\textbf{Obstetricia}   & Control prenatal, biometría fetal, Doppler umbilical --- sin radiación \\[2pt]
\textbf{Cardiología}   & Ecocardiografía transtorácica y transesofágica --- función ventricular \\[2pt]
\textbf{Vascular}      & Doppler carotídeo, venoso profundo --- diagnóstico TVP y estenosis \\[2pt]
\textbf{MSK/Urgencias} & Fracturas, derrames articulares, guía de procedimientos en tiempo real \\
\end{tabular}
\end{clinicbox}

\newpage

%=============================================================================
\section{Radiografía Convencional y Digital (RX --- X-Ray)}
%=============================================================================

\begin{principiobox}
Atenuación diferencial de rayos X según densidad y número atómico ($Z$) del tejido.
Los sistemas digitales DR (panel plano FPD) reemplazan al CR con mayor DQE y menor latencia.
Imagen 2D proyeccional.
\end{principiobox}

\subsection{Comparativa de parámetros técnicos}

\begin{center}
\renewcommand{\arraystretch}{1.25}
\begin{tabular}{P{4.2cm} C{2.5cm} C{2.5cm} C{2.5cm} C{2.5cm}}
\toprule
\rowcolor{navy}
\textbf{\color{white}Parámetro} &
\textbf{\color{white}RX Conv. (CR)} &
\textbf{\color{white}RX Digital (DR)} &
\textbf{\color{white}RX Móvil DR} &
\textbf{\color{white}Mamógrafo Dig.} \\
\midrule
Detector                   & Fósforo (CR) & \cellcolor{bestgreen}$\bigstar$ FPD Si amorfo & FPD portátil & FPD Se amorfo \\
\rowcolor{rowalt}
Resolución (lp/mm)         & 3.0--4.0 & 3.4--4.5 & 3.0--3.5 & \cellcolor{bestgreen}$\bigstar$ 5.0--7.0 \\
Rango kVp                  & 40--150 & 40--150 & 40--125 & 25--35 \\
\rowcolor{rowalt}
Corriente máx. (mA)        & 500 & \cellcolor{bestgreen}$\bigstar$ 800 & 300 & 100 \\
Latencia imagen (s)        & 60--120 & \cellcolor{bestgreen}$\bigstar$ 3--5 & \cellcolor{bestgreen}$\bigstar$ 3--5 & 5--10 \\
\rowcolor{rowalt}
Tamaño detector (cm)       & 35$\times$43 & \cellcolor{bestgreen}$\bigstar$ 43$\times$43 & 35$\times$43 & 24$\times$30 \\
DQE (\%)                   & 30--45 & \cellcolor{bestgreen}$\bigstar$ 65--75 & 55--65 & \cellcolor{bestgreen}$\bigstar$ 70--80 \\
\rowcolor{rowalt}
DICOM compatible           & Sí & Sí & Sí & Sí \\
Tomosíntesis               & No & Opcional & No & \cellcolor{bestgreen}$\bigstar$ Sí \\
\rowcolor{rowalt}
IA para lectura            & No & \cellcolor{bestgreen}$\bigstar$ Sí & Básica & \cellcolor{bestgreen}$\bigstar$ Sí \\
Precio ref. (USD)          & 15K--40K & 80K--180K & 40K--90K & 150K--300K \\
\bottomrule
\end{tabular}
\end{center}

\begin{itemize}[leftmargin=*, itemsep=2pt, topsep=4pt]
  \footnotesize
  \item \textbf{DQE} (\textit{Detective Quantum Efficiency}): eficiencia de detección de fotones --- mayor valor = mejor relación calidad/dosis
  \item \textbf{Tomosíntesis (DBT)}: múltiples proyecciones para reconstrucción pseudo-3D --- reduce solapamiento tisular
  \item \textbf{RX móvil}: UCI, quirófano, neonatal --- sin necesidad de trasladar al paciente
\end{itemize}

\begin{clinicbox}
\begin{tabular}{@{}P{3cm} P{10.5cm}@{}}
\textbf{Tórax}          & Primera línea: neumonía, derrame, neumotórax, cardiomegalia \\[2pt]
\textbf{Traumatología}  & Fracturas, luxaciones, cuerpos extraños --- alta disponibilidad \\[2pt]
\textbf{Columna}        & Evaluación postural, escoliosis, espondilopatías \\[2pt]
\textbf{Abdomen}        & Obstrucción intestinal, neumoperitoneo --- proyección en bipedestación \\[2pt]
\textbf{Mamografía}     & Tamizaje y diagnóstico de cáncer de mama --- protocolo ACR \\
\end{tabular}
\end{clinicbox}

\newpage

%=============================================================================
\section{PET/CT --- Tomografía por Emisión de Positrones}
%=============================================================================

\begin{principiobox}
Detección en coincidencia de pares de fotones gamma (511~keV) emitidos por aniquilación
positrón--electrón. El radiotrazador más utilizado es $^{18}$F-FDG (fluorodesoxiglucosa).
La fusión PET/CT correlaciona actividad metabólica con anatomía.
\end{principiobox}

\vspace{6pt}
El tiempo de vuelo (\textbf{TOF}) mejora la resolución de imagen al calcular la diferencia
temporal entre la detección de los dos fotones gamma, localizando la emisión en el eje axial
con mayor precisión. Los sistemas digitales de tercera generación alcanzan resoluciones
espaciales de 2.5--3~mm y sensibilidades de 15--20~cps/kBq.

\subsection{Comparativa de parámetros técnicos}

\begin{center}
\renewcommand{\arraystretch}{1.25}
\begin{tabular}{P{4.5cm} C{2.3cm} C{2.3cm} C{2.3cm} C{2.3cm}}
\toprule
\rowcolor{navy}
\textbf{\color{white}Parámetro} &
\textbf{\color{white}PET/CT Estándar} &
\textbf{\color{white}PET/CT TOF} &
\textbf{\color{white}PET/CT Digital} &
\textbf{\color{white}PET/MR Híbrido} \\
\midrule
Resolución espacial (mm)   & 4--5 & 3--4 & \cellcolor{bestgreen}$\bigstar$ 2.5--3 & \cellcolor{bestgreen}$\bigstar$ 2.5--3 \\
\rowcolor{rowalt}
Sensibilidad (cps/kBq)     & 5--8 & 8--12 & \cellcolor{bestgreen}$\bigstar$ 12--18 & \cellcolor{bestgreen}$\bigstar$ 15--20 \\
TOF (ps)                   & No & 400--600 & \cellcolor{bestgreen}$\bigstar$ 210--370 & \cellcolor{bestgreen}$\bigstar$ 230--400 \\
\rowcolor{rowalt}
Longitud axial (cm)        & 15--22 & 22--26 & \cellcolor{bestgreen}$\bigstar$ 26--32 & 25--30 \\
Componente CT (cortes)     & 16--64 & \cellcolor{bestgreen}$\bigstar$ 64--128 & \cellcolor{bestgreen}$\bigstar$ 128--256 & No aplica \\
\rowcolor{rowalt}
Componente MR              & No & No & No & \cellcolor{bestgreen}$\bigstar$ Sí (3T) \\
Radiotrazadores            & FDG & FDG, Ga-68 & \cellcolor{bestgreen}$\bigstar$ Múltiples & \cellcolor{bestgreen}$\bigstar$ Múltiples \\
\rowcolor{rowalt}
Dosis radiación (mSv)      & \cellcolor{bestgreen}$\bigstar$ 3--5 & \cellcolor{bestgreen}$\bigstar$ 3--5 & \cellcolor{bestgreen}$\bigstar$ 2--4 & \cellcolor{bestgreen}$\bigstar$ 1--3 \\
Tiempo adquisición (min)   & 20--30 & 15--25 & \cellcolor{bestgreen}$\bigstar$ 10--20 & 45--90 \\
\rowcolor{rowalt}
IA en reconstrucción       & No & Básica & \cellcolor{bestgreen}$\bigstar$ Avanzada & \cellcolor{bestgreen}$\bigstar$ Avanzada \\
Precio ref. (USD M)        & 1.5--2.5 & 2.0--3.0 & 2.5--4.0 & 5.0--8.0 \\
\bottomrule
\end{tabular}
\end{center}

\begin{itemize}[leftmargin=*, itemsep=2pt, topsep=4pt]
  \footnotesize
  \item \textbf{TOF} (\textit{Time of Flight}): mejora SNR localizando la emisión mediante diferencia temporal entre fotones
  \item \textbf{$^{18}$F-FDG}: detecta hipermetabolismo en tejido tumoral y cerebral activo
  \item \textbf{PET/MR}: máxima información combinada --- uso principalmente en investigación y oncología pediátrica
\end{itemize}

\begin{clinicbox}
\begin{tabular}{@{}P{3.5cm} P{10cm}@{}}
\textbf{Oncología}          & Estadificación, respuesta a tratamiento, detección de recidiva \\[2pt]
\textbf{Neurología}         & Diagnóstico diferencial demencias (Alzheimer, FTD), epilepsia refractaria \\[2pt]
\textbf{Cardiología}        & Viabilidad miocárdica, sarcoidosis cardíaca, endocarditis \\[2pt]
\textbf{Infecciones}        & Fiebre de origen desconocido, osteomielitis, vasculitis de grandes vasos \\[2pt]
\textbf{Planificación RT}   & Definición de volumen tumoral para radioterapia de precisión \\
\end{tabular}
\end{clinicbox}

\newpage

%=============================================================================
\section{Fluoroscopía y Sistemas de Rayos X Intervencionistas}
%=============================================================================

\begin{principiobox}
Imagen de rayos X en tiempo real mediante detector de panel plano (FPD) o intensificador
de imagen clásico. Exposición continua o pulsada. Permite guía en tiempo real de
procedimientos intervencionistas mínimamente invasivos.
\end{principiobox}

\subsection{Comparativa de parámetros técnicos}

\begin{center}
\renewcommand{\arraystretch}{1.25}
\begin{tabular}{P{4.2cm} C{2.5cm} C{2.5cm} C{2.5cm} C{2.5cm}}
\toprule
\rowcolor{navy}
\textbf{\color{white}Parámetro} &
\textbf{\color{white}Fluoro General} &
\textbf{\color{white}Arco en C Móvil} &
\textbf{\color{white}Sala Hemo.} &
\textbf{\color{white}Sala Híbrida} \\
\midrule
Detector                   & FPD 43$\times$43~cm & FPD 30$\times$30~cm & \cellcolor{bestgreen}$\bigstar$ FPD 30$\times$40~cm & \cellcolor{bestgreen}$\bigstar$ FPD 30$\times$40~cm \\
\rowcolor{rowalt}
Resolución (lp/mm)         & 2.5--3.5 & 2.0--2.8 & \cellcolor{bestgreen}$\bigstar$ 3.0--4.0 & \cellcolor{bestgreen}$\bigstar$ 3.5--4.5 \\
Rango kVp                  & 40--125 & 40--120 & \cellcolor{bestgreen}$\bigstar$ 40--125 & \cellcolor{bestgreen}$\bigstar$ 40--125 \\
\rowcolor{rowalt}
Dosis fluoro (mGy/min)     & \cellcolor{bestgreen}$\bigstar$ 1--5 & 2--8 & \cellcolor{bestgreen}$\bigstar$ 1--4 & \cellcolor{bestgreen}$\bigstar$ 1--3 \\
DAP meter integrado        & Sí & Sí & \cellcolor{bestgreen}$\bigstar$ Avanzado & \cellcolor{bestgreen}$\bigstar$ Avanzado \\
\rowcolor{rowalt}
DSA                        & Opcional & No & \cellcolor{bestgreen}$\bigstar$ Sí & \cellcolor{bestgreen}$\bigstar$ Sí \\
Rotación isocentro (°)     & 0 (fijo) & $\pm$180 & \cellcolor{bestgreen}$\bigstar$ $\pm$270 & \cellcolor{bestgreen}$\bigstar$ $\pm$270 \\
\rowcolor{rowalt}
Angiografía 3D             & No & No & \cellcolor{bestgreen}$\bigstar$ Sí & \cellcolor{bestgreen}$\bigstar$ Sí \\
Integración TC/MR          & No & No & No & \cellcolor{bestgreen}$\bigstar$ Sí \\
\rowcolor{rowalt}
IA guidance                & No & No & Sí & \cellcolor{bestgreen}$\bigstar$ Avanzada \\
Precio ref. (USD)          & 80K--200K & 40K--120K & 800K--2.0M & 2.0M--5.0M \\
\bottomrule
\end{tabular}
\end{center}

\begin{itemize}[leftmargin=*, itemsep=2pt, topsep=4pt]
  \footnotesize
  \item \textbf{DAP} (\textit{Dose-Area Product}): métrica dosimétrica estándar en fluoroscopía --- requisito regulatorio
  \item \textbf{DSA} (\textit{Digital Subtraction Angiography}): sustracción de imagen basal para visualizar contraste vascular
  \item \textbf{Sala híbrida}: combina sala quirúrgica estéril con arco de alta calidad --- cirugía cardiovascular y endovascular
\end{itemize}

\begin{clinicbox}
\begin{tabular}{@{}P{3.5cm} P{10cm}@{}}
\textbf{Cardiología Interv.}  & Cateterismo, angioplastia, TAVI, implante de stents coronarios \\[2pt]
\textbf{Radiología Interv.}   & Embolización, drenajes, biopsias guiadas, TIPS \\[2pt]
\textbf{Ortopedia}            & Reducción de fracturas, osteosíntesis, vertebroplastia \\[2pt]
\textbf{Gastroenterología}    & CPRE, colocación de stents biliares, manometría \\[2pt]
\textbf{Neurología Interv.}   & Angiografía cerebral, embolización de aneurismas, trombectomía \\
\end{tabular}
\end{clinicbox}

\newpage

%=============================================================================
\section*{Notas sobre especificaciones técnicas}
\addcontentsline{toc}{section}{Notas sobre especificaciones técnicas}
%=============================================================================

\begin{itemize}[leftmargin=*, itemsep=5pt]
  \item Los valores presentados provienen de \textit{datasheets} públicos de fabricantes
        (2024--2025) y pueden variar según configuración, accesorios opcionales y mercado.
  \item Los rangos de precio son \textbf{referenciales FOB} y no incluyen instalación,
        obras civiles, capacitación ni contrato de mantenimiento.
  \item Los fabricantes se identifican genéricamente (\textit{Marca A, B, C, D}) para
        mantener neutralidad en el análisis comparativo.
  \item La celda $\bigstar$ indica el valor más favorable dentro de esa categoría de
        parámetro.
\end{itemize}

\vspace{8pt}

\noindent\textbf{Referencias y fuentes consultadas}

\vspace{4pt}
\begin{itemize}[leftmargin=*, itemsep=3pt]
  \item Bushberg, J.T. et al. \textit{The Essential Physics of Medical Imaging}, 3rd ed.
        Lippincott Williams \& Wilkins, 2011.
  \item Sprawls, P. \textit{Physical Principles of Medical Imaging}, 2nd ed.
        Medical Physics Publishing, 1995.
  \item IEC 60601-2-33: \textit{Particular requirements for the basic safety and essential
        performance of magnetic resonance equipment for medical diagnosis.}
  \item WHO. \textit{Medical Devices: Managing the Mismatch.} WHO Press, 2010.
  \item ECRI Institute. \textit{Medical Equipment Comparison Reports} (2024--2025).
  \item Portales de producto: sitios públicos de los principales fabricantes de
        equipos de imagen médica (consultados 2024--2025).
\end{itemize}

\vspace{16pt}
\noindent{\color{graybrd}\rule{\linewidth}{0.5pt}}
\vspace{6pt}

\noindent\small\color{graym}
\textit{Joselyn Romero Avila · B.Sc. Ingeniería Biomédica · UNMSM · 2025}

\end{document}
